\chapter{Polynomials and the FFT}

Here is a small mystery to chew on before we begin. Multiplying two big numbers,
or two polynomials, feels like it should cost a fixed amount of work --- you learned
the ``multiply every digit by every digit'' method in grade school, and it never
occurred to you that it might be slow. But it \emph{is} slow. To multiply two
polynomials of degree $n$ the schoolbook method does about $n^2$ little
multiplications. For $n = 1{,}000{,}000$ that is a trillion operations. For a long
time people simply assumed there was nothing to be done; multiplication was
multiplication.

This chapter is the story of how that assumption was wrong, and of the beautiful
divide-and-conquer algorithm --- the \emph{Fast Fourier Transform}, or FFT --- that
breaks $n^2$ down to $n\log n$. The same trick that makes the FFT fast is the trick
that makes your phone decode an MP3, that lets a radio telescope clean up an image,
and that underlies a large part of digital signal processing. We are going to build
it from the ground up, starting with something completely innocent: two different
ways to write down a polynomial.

\begin{focusbox}
By the end of this chapter you should be able to: (1) describe a polynomial in
\emph{coefficient} form and in \emph{point-value} form, and say why each is good at
a different job; (2) explain why multiplying in point-value form takes only linear
time; (3) define the roots of unity and list the three algebraic facts about them
that make the FFT work; (4) derive the FFT as a divide-and-conquer algorithm and
prove its $\Th(n\log n)$ running time from a recurrence; and (5) assemble the whole
pipeline that multiplies two polynomials in $\Th(n\log n)$ time.
\end{focusbox}

\section{Two ways to describe a polynomial}

A polynomial of degree bound $n$ (meaning degree at most $n-1$) is the object
$$
A(x) = a_0 + a_1 x + a_2 x^2 + \cdots + a_{n-1}x^{n-1} = \sum_{j=0}^{n-1} a_j x^j.
$$
The numbers $a_0, a_1, \dots, a_{n-1}$ are its \emph{coefficients}. This is the
representation you have used your whole life, and it is so familiar that you may
never have noticed it is a \emph{choice}. We will call the list
$(a_0, a_1, \dots, a_{n-1})$ the \emph{coefficient representation} of $A$.

\begin{intuitionbox}
A polynomial is a function: feed it a number $x$, get a number $A(x)$ back. The
coefficient list is one way to pin down which function it is --- it is like
describing a person by their DNA. But there is a totally different way to pin down
the same function: instead of listing its genetic recipe, you list a bunch of
\emph{things it actually does}. ``At $x=0$ it outputs $3$; at $x=1$ it outputs $7$;
at $x=2$ it outputs $17$.'' If you give enough such observations, they nail down the
function just as tightly as the coefficients do.
\end{intuitionbox}

\subsection{Point-value form}

Here is the second representation. A \emph{point-value representation} of a
polynomial $A$ of degree bound $n$ is a set of $n$ point-value pairs
$$
\set{(x_0, y_0), (x_1, y_1), \dots, (x_{n-1}, y_{n-1})}
$$
where the $x_k$ are all distinct and $y_k = A(x_k)$ for each $k$. So instead of the
recipe, you record what the polynomial \emph{evaluates to} at $n$ chosen inputs.

Why exactly $n$ points, no more and no fewer? Because of a fact you may remember
from school dressed up in a fancier outfit: through any two points there is exactly
one line ($\deg \le 1$, so $2$ coefficients); through any three points in general
position there is exactly one parabola ($\deg \le 2$, so $3$ coefficients); and in
general, $n$ points determine a unique polynomial of degree bound $n$.

\begin{theorem}[Uniqueness of interpolation]\label{ch29:interp-thm}
For any set of $n$ point-value pairs $\set{(x_0,y_0),\dots,(x_{n-1},y_{n-1})}$ with
all $x_k$ distinct, there is a \emph{unique} polynomial $A$ of degree bound $n$ such
that $A(x_k)=y_k$ for all $k$.
\end{theorem}

\begin{formalbox}
\textbf{Why it is true.} Asking for $A(x_k)=y_k$ for $k=0,\dots,n-1$ is asking the
unknown coefficients $a_0,\dots,a_{n-1}$ to satisfy the linear system
$$
\begin{pmatrix}
1 & x_0 & x_0^2 & \cdots & x_0^{n-1}\\
1 & x_1 & x_1^2 & \cdots & x_1^{n-1}\\
\vdots & & & & \vdots\\
1 & x_{n-1} & x_{n-1}^2 & \cdots & x_{n-1}^{n-1}
\end{pmatrix}
\begin{pmatrix}a_0\\a_1\\\vdots\\a_{n-1}\end{pmatrix}
=
\begin{pmatrix}y_0\\y_1\\\vdots\\y_{n-1}\end{pmatrix}.
$$
The matrix on the left is the \emph{Vandermonde matrix} $V(x_0,\dots,x_{n-1})$. Its
determinant equals $\prod_{j<k}(x_k - x_j)$, which is nonzero exactly when all the
$x_k$ are distinct. A square linear system with a nonzero determinant has exactly
one solution --- so the coefficients exist and are unique. That is the whole proof.
\end{formalbox}

The two representations are therefore \emph{equivalent}: each one determines the
same polynomial, and we may convert back and forth.

\begin{itemize}
\item \textbf{Evaluation} goes coefficient $\to$ point-value: pick the points $x_k$,
compute $y_k=A(x_k)$.
\item \textbf{Interpolation} goes point-value $\to$ coefficient: recover the $a_j$
from the pairs $(x_k,y_k)$. (This is solving the Vandermonde system, or running
Lagrange's formula.)
\end{itemize}

\begin{focusbox}
Hold these two pictures side by side for the whole chapter:

\textbf{coefficient form} $(a_0,\dots,a_{n-1})$ \quad$\xrightleftharpoons[\text{interpolate}]{\text{evaluate}}$\quad \textbf{point-value form} $\set{(x_k,y_k)}$.

The FFT is nothing more than a \emph{breathtakingly fast} way to travel along this
double arrow when the $x_k$ are chosen cleverly.
\end{focusbox}

\subsection{A small worked conversion}

Take $A(x) = 1 + 2x + 3x^2$ (so $n=3$, coefficients $(1,2,3)$). Let us evaluate at
$x_0=0$, $x_1=1$, $x_2=2$.
$$
A(0) = 1,\qquad
A(1) = 1+2+3 = 6,\qquad
A(2) = 1+4+12 = 17.
$$
So a point-value representation is $\set{(0,1),(1,6),(2,17)}$. Three coefficients in,
three pairs out --- both forms carry exactly the same information about the same
parabola. If somebody hands you those three pairs and you did not know the
coefficients, Theorem~\ref{ch29:interp-thm} promises you could recover $(1,2,3)$ and
nothing else.

\begin{center}
\begin{tikzpicture}[>=Stealth,xscale=1.6,yscale=0.18]
\draw[->,gray] (-0.4,0)--(2.6,0) node[right,font=\scriptsize]{$x$};
\draw[->,gray] (0,-1)--(0,19) node[above,font=\scriptsize]{$A(x)$};
\foreach \x in {0,1,2} \draw[gray!40] (\x,-0.6)--(\x,0.6) node[below=2pt,font=\tiny] {$\x$};
% parabola A(x)=1+2x+3x^2
\draw[sourceblue,thick,domain=-0.3:2.25,samples=60,smooth]
  plot (\x,{1+2*\x+3*\x*\x});
\foreach \x/\y in {0/1,1/6,2/17}{
  \node[circle,fill=structurepurple,inner sep=1.5pt] at (\x,\y) {};
  \node[font=\tiny,right=1pt] at (\x,\y) {$(\x,\y)$};
}
\end{tikzpicture}
\end{center}

The three purple dots are the point-value pairs; the blue curve is the unique
parabola threaded through them. Move any dot and you get a different parabola ---
but fix all three and the curve has no freedom left, which is exactly what
Theorem~\ref{ch29:interp-thm} asserts.

\section{The operation we actually care about: multiplication}

Now let us see why anyone would bother with two representations. Suppose we want the
\emph{product} of two polynomials,
$$
A(x) = \sum_{j=0}^{n-1} a_j x^j,\qquad B(x) = \sum_{j=0}^{n-1} b_j x^j,
$$
giving a polynomial $C(x) = A(x)B(x)$ of degree bound $2n-1$. In coefficient form,
its coefficients are the \emph{convolution} of the two coefficient lists:
$$
c_k = \sum_{j=0}^{k} a_j\, b_{k-j}, \qquad k = 0,1,\dots,2n-2.
$$

\begin{intuitionbox}
Convolution is just ``collect every way to make $x^k$.'' To get the coefficient of
$x^k$ in the product, you pair up a term $a_j x^j$ from $A$ with the term
$b_{k-j}x^{k-j}$ from $B$, because their exponents add to $k$. You sum over all such
pairs. Computing one $c_k$ costs about $k$ multiplications; summing over all $k$
costs $\Th(n^2)$. That is the schoolbook cost, and it is exactly what we want to
beat.
\end{intuitionbox}

Now watch what happens in \emph{point-value} form. Suppose we already know $A$ and
$B$ at the \emph{same} set of points $x_0,\dots,x_{m-1}$:
$$
A: \set{(x_k, A(x_k))},\qquad B:\set{(x_k, B(x_k))}.
$$
Then the value of the product at any one point is just the product of the two
values, because $C(x_k) = A(x_k)\,B(x_k)$ --- multiplication of \emph{functions} is
multiplication of their \emph{outputs}, pointwise. So
$$
C: \set{(x_k,\; A(x_k)\,B(x_k))}.
$$

\begin{focusbox}
\textbf{Point-value multiplication is embarrassingly cheap.} If both polynomials are
sampled at the same points, multiplying them is just multiplying the $y$-values one
at a time: $m$ scalar multiplications, total cost $\Th(m)$. Linear! Compare that
to the $\Th(n^2)$ convolution in coefficient form.
\end{focusbox}

There is one bookkeeping subtlety. The product $C$ has degree bound $2n-1$, so by
Theorem~\ref{ch29:interp-thm} it needs at least $2n-1$ points to be pinned down. So
before sampling we treat $A$ and $B$ as degree-bound-$2n$ polynomials by padding their
coefficient lists with zeros, and we evaluate both at $2n$ points. Then pointwise
multiplication gives a valid point-value form of $C$.

This suggests a strategy with three steps, drawn here as a journey:

\begin{center}
\begin{tikzpicture}[node distance=10mm and 16mm,
  box/.style={draw,rounded corners,align=center,inner sep=4pt,font=\small},
  >=Stealth]
\node[box,fill=blue!4] (coefAB) {coefficients of $A,B$\\(padded to length $2n$)};
\node[box,fill=green!5,right=of coefAB] (pvAB) {values of $A,B$\\at $2n$ points};
\node[box,fill=green!5,right=of pvAB] (pvC) {values of $C$\\at the same points};
\node[box,fill=blue!4,below=14mm of pvC] (coefC) {coefficients of $C$};
\draw[->,sourceblue,thick] (coefAB) -- node[above,font=\scriptsize]{\textbf{evaluate}} node[below,font=\scriptsize]{(2 FFTs)} (pvAB);
\draw[->,targetgreen,thick] (pvAB) -- node[above,align=center,font=\scriptsize]{\textbf{pointwise}\\ \textbf{multiply}} node[below,font=\scriptsize]{$\Th(n)$} (pvC);
\draw[->,quotientorange,thick] (pvC) -- node[right,align=left,font=\scriptsize]{\textbf{interpolate}\\(inverse FFT)} (coefC);
\end{tikzpicture}
\end{center}

The middle step is already linear. The whole battle is in the side steps:
\emph{evaluate} (go right) and \emph{interpolate} (come back). Done naively, each of
those is $\Th(n^2)$ --- evaluating one polynomial at $n$ points by Horner's rule is
$\Th(n)$ per point, $\Th(n^2)$ for all points --- and we have gained nothing. The
FFT's entire job is to make those two side trips cost $\Th(n\log n)$. And the secret
is to stop letting the points $x_k$ be arbitrary, and instead choose them with
exquisite care.

\section{Choosing the points: roots of unity}

Why should the choice of evaluation points matter for \emph{speed}? Here is the
intuition, and it is the soul of the whole algorithm.

\begin{intuitionbox}
Divide-and-conquer wants to split a problem of size $n$ into two problems of size
$n/2$ that we can solve and cheaply recombine. To evaluate $A$ at $n$ points, we
would love to split $A$ into two smaller polynomials and evaluate \emph{those} at
$n/2$ points each. For that to save work, the $n/2$ smaller points must somehow be
``shared'' between the two halves --- evaluating the big polynomial at a point and
its negative should reuse the same sub-computation. Ordinary points like
$0,1,2,3,\dots$ have no such hidden relationships. But there is a magical set of
points, evenly spaced around a circle, where squaring the set gives back a set half
its size. Those are the \emph{roots of unity}.
\end{intuitionbox}

\subsection{Complex roots of unity}

A \emph{complex $n$th root of unity} is a number $\omega$ with $\omega^n = 1$. There
are exactly $n$ of them, and they sit at equal angles around the unit circle in the
complex plane. Using Euler's formula $e^{i\theta} = \cos\theta + i\sin\theta$, the
\emph{principal} $n$th root of unity is
$$
\omega_n = e^{2\pi i/n},
$$
and the full set of $n$ roots is $\omega_n^0, \omega_n^1, \dots, \omega_n^{n-1}$.
They are the $n$ points you reach by walking around the unit circle in $n$ equal
steps of angle $2\pi/n$.

\begin{center}
\begin{tikzpicture}[scale=1.5,>=Stealth]
\draw[gray!50] (0,0) circle (1);
\draw[->,gray] (-1.25,0)--(1.3,0) node[right,font=\scriptsize]{Re};
\draw[->,gray] (0,-1.25)--(0,1.3) node[above,font=\scriptsize]{Im};
\foreach \k/\lab in {0/{\omega_8^0=1}, 1/{\omega_8^1}, 2/{\omega_8^2=i}, 3/{\omega_8^3}, 4/{\omega_8^4=-1}, 5/{\omega_8^5}, 6/{\omega_8^6=-i}, 7/{\omega_8^7}}{
  \node[circle,fill=structurepurple,inner sep=1.4pt] (p\k) at ({cos(\k*45)},{sin(\k*45)}) {};
}
\node[right=1pt,font=\scriptsize] at (p0) {$\omega_8^0=1$};
\node[above right=0pt,font=\scriptsize] at (p1) {$\omega_8^1$};
\node[above,font=\scriptsize] at (p2) {$\omega_8^2=i$};
\node[above left=0pt,font=\scriptsize] at (p3) {$\omega_8^3$};
\node[left=1pt,font=\scriptsize] at (p4) {$\omega_8^4=-1$};
\node[below left=0pt,font=\scriptsize] at (p5) {$\omega_8^5$};
\node[below,font=\scriptsize] at (p6) {$\omega_8^6=-i$};
\node[below right=0pt,font=\scriptsize] at (p7) {$\omega_8^7$};
\end{tikzpicture}
\end{center}

The figure shows the eight $8$th roots of unity. Notice they come in
\emph{antipodal pairs}: $\omega_8^0$ and $\omega_8^4=-1$ are opposite; so are
$\omega_8^1$ and $\omega_8^5$; and so on. That pairing is not a coincidence, and it
is precisely the structure divide-and-conquer will exploit.

\subsection{Three facts that make everything work}

Everything the FFT needs rests on three small algebraic facts. Let us state and feel
each one. Assume $n$ is even (in fact a power of two).

\begin{lemma}[Cancellation]\label{ch29:cancel}
For any integers $n>0$, $d>0$, and $k\ge 0$, $\ \omega_{dn}^{dk} = \omega_n^{k}.$
In particular $\omega_n^{k+n/2} = -\omega_n^{k}$ (taking $d=2$ together with
$\omega_2^1 = -1$).
\end{lemma}

\begin{formalbox}
\textbf{Why.} $\omega_{dn}^{dk} = \big(e^{2\pi i/(dn)}\big)^{dk}
= e^{2\pi i \, dk/(dn)} = e^{2\pi i\, k/n} = \omega_n^k.$ The exponents just cancel.
For the second statement, $\omega_n^{k+n/2} = \omega_n^k\,\omega_n^{n/2}$, and
$\omega_n^{n/2} = e^{2\pi i \,(n/2)/n} = e^{\pi i} = -1$. So evaluating at a root and
at its antipode differ only by a sign --- this is the ``share the work between a
point and its negative'' phenomenon we hoped for.
\end{formalbox}

\begin{lemma}[Halving]\label{ch29:halving}
If $n>0$ is even, then the squares of the $n$ complex $n$th roots of unity are the
$n/2$ complex $(n/2)$th roots of unity, \emph{each appearing exactly twice}.
\end{lemma}

\begin{formalbox}
\textbf{Why.} Square the $k$th root: $\big(\omega_n^k\big)^2 = \omega_n^{2k}
= \omega_{n/2}^{k}$ by cancellation (with $d=2$, replacing $n$ by $n/2$). As $k$
runs $0,1,\dots,n-1$, the value $\omega_{n/2}^{k}$ cycles through all $n/2$ of the
$(n/2)$th roots, and because $\omega_n^k$ and $\omega_n^{k+n/2}$ square to the same
thing (their squares differ by $\omega_n^{n} = 1$), each value is hit exactly twice.
\end{formalbox}

\begin{intuitionbox}
The halving lemma is the entire reason the FFT exists. It says: if I want to
evaluate a polynomial at all $n$ roots of unity, the \emph{squares} of those points
form a problem half the size --- the $n/2$ roots of unity --- and the squares are
what a divide step naturally produces. So a size-$n$ evaluation reduces to a
size-$n/2$ evaluation. Halve, halve, halve: that is where the $\log n$ comes from.
\end{intuitionbox}

\begin{lemma}[Summation]\label{ch29:summation}
For any integer $k$ not divisible by $n$,
$\displaystyle \sum_{j=0}^{n-1}\big(\omega_n^{k}\big)^{j} = 0.$
(When $k$ \emph{is} divisible by $n$, every term is $1$ and the sum is $n$.)
\end{lemma}

\begin{formalbox}
\textbf{Why.} This is a geometric series with ratio $\omega_n^k$:
$$
\sum_{j=0}^{n-1}(\omega_n^k)^j
= \frac{(\omega_n^k)^n - 1}{\omega_n^k - 1}
= \frac{(\omega_n^n)^k - 1}{\omega_n^k - 1}
= \frac{1^k - 1}{\omega_n^k - 1} = 0,
$$
valid because $\omega_n^k \ne 1$ when $n\nmid k$, so we are not dividing by zero. The
summation lemma is what makes \emph{interpolation} (the trip back) work; we will use
it when we invert the FFT.
\end{formalbox}

\section{The Discrete Fourier Transform}

We are now ready to name the object we are computing. We evaluate $A$ at the $n$
roots of unity. The resulting list of values is so important it gets its own name.

\begin{definition}[DFT]\label{ch29:dft-def}
Let $A$ have coefficient vector $(a_0,\dots,a_{n-1})$. Its \emph{Discrete Fourier
Transform} is the vector $(y_0,\dots,y_{n-1})$ where
$$
y_k = A(\omega_n^k) = \sum_{j=0}^{n-1} a_j \,\omega_n^{kj},\qquad k=0,1,\dots,n-1.
$$
We write $\mathbf{y} = \mathrm{DFT}_n(\mathbf{a})$.
\end{definition}

So the DFT is just ``evaluate at the roots of unity,'' packaged. Computed directly
from the definition --- $n$ values, each a sum of $n$ terms --- it costs $\Th(n^2)$.
The \emph{Fast} Fourier Transform is an algorithm that produces the very same vector
$\mathbf{y}$ in $\Th(n\log n)$ time. Same output; cleverer route.

\begin{aestheticbox}
The DFT is one of the most used computations on Earth. Sampled sound, the pixels
streaming off a camera sensor, the radio waves hitting a cellphone tower, the
seismograph trace of an earthquake --- all of these are sequences of numbers, and
asking ``which frequencies are present, and how strong?'' is exactly asking for
their DFT. Without a \emph{fast} way to compute it, digital audio, image
compression (JPEG), wireless communication (OFDM in Wi-Fi and 4G/5G), and medical
imaging would all be impractical. Gilbert Strang has called the FFT ``the most
important numerical algorithm of our lifetime.''
\end{aestheticbox}

\section{The FFT: divide and conquer on the coefficients}

Here comes the idea. We want $A(\omega_n^k)$ for all $k$, where
$A(x) = a_0 + a_1 x + \cdots + a_{n-1}x^{n-1}$ and $n$ is a power of two. Split $A$
into its \emph{even-indexed} and \emph{odd-indexed} coefficients:
$$
A^{[0]}(x) = a_0 + a_2 x + a_4 x^2 + \cdots + a_{n-2}x^{n/2 - 1},
$$
$$
A^{[1]}(x) = a_1 + a_3 x + a_5 x^2 + \cdots + a_{n-1}x^{n/2 - 1}.
$$
Each of these is a polynomial of degree bound $n/2$ --- half the size. And they
reassemble the original by
$$
A(x) = A^{[0]}(x^2) + x\,A^{[1]}(x^2).
$$

\begin{intuitionbox}
Read that identity slowly. Group the even-power terms: $a_0 + a_2x^2 + a_4x^4 +\cdots$
is exactly $A^{[0]}$ evaluated at $x^2$. Group the odd-power terms:
$a_1x + a_3x^3+\cdots = x(a_1 + a_3 x^2 + \cdots) = x\,A^{[1]}(x^2)$. Add them, you
recover $A(x)$. So to evaluate $A$ at a point $x$, it suffices to evaluate the two
half-size polynomials at $x^2$ and combine with one multiply and one add.
\end{intuitionbox}

Now plug in the roots of unity for $x$. To get all of $A(\omega_n^0),\dots,
A(\omega_n^{n-1})$, we need $A^{[0]}$ and $A^{[1]}$ evaluated at the \emph{squares}
$\big(\omega_n^k\big)^2$. But by the halving lemma those squares are exactly the
$n/2$ roots $\omega_{n/2}^{k}$ --- and there are only $n/2$ distinct ones! So:

\begin{focusbox}
Evaluating $A$ at the $n$ roots of unity reduces to evaluating $A^{[0]}$ and
$A^{[1]}$ --- each half the degree --- at the $n/2$ roots of unity. Two
subproblems of half the size. That is a clean divide-and-conquer split, and the
halving lemma is what guarantees the recursion lands back on the same kind of
problem.
\end{focusbox}

\subsection{The combine step, made explicit}

Suppose the two recursive calls have returned. Let
$y^{[0]}_k = A^{[0]}(\omega_{n/2}^{k})$ and $y^{[1]}_k = A^{[1]}(\omega_{n/2}^{k})$
for $k=0,\dots,n/2-1$. We want $y_k = A(\omega_n^k)$ for $k = 0,\dots,n-1$.

For $k$ in the \emph{lower} half, $0 \le k < n/2$, plug $x=\omega_n^k$ into the
identity and use $\big(\omega_n^k\big)^2 = \omega_{n/2}^k$:
$$
y_k = A^{[0]}(\omega_{n/2}^k) + \omega_n^k\, A^{[1]}(\omega_{n/2}^k)
    = y^{[0]}_k + \omega_n^k\, y^{[1]}_k.
$$
For the \emph{upper} half, look at index $k+n/2$. Here is where the cancellation
lemma pays off. We have $\big(\omega_n^{k+n/2}\big)^2 = \omega_{n/2}^{k}$ (squaring
removes the antipodal sign), so $A^{[0]}$ and $A^{[1]}$ are evaluated at the
\emph{same} point as before; and $\omega_n^{k+n/2} = -\omega_n^k$. Therefore
$$
y_{k+n/2} = A^{[0]}(\omega_{n/2}^k) - \omega_n^k\, A^{[1]}(\omega_{n/2}^k)
          = y^{[0]}_k - \omega_n^k\, y^{[1]}_k.
$$

\begin{intuitionbox}
Two outputs for the price of one. The quantity $\omega_n^k\, y^{[1]}_k$ --- called
the \emph{twiddle factor} times the odd value --- is computed once, then \emph{added}
to give the lower-half answer and \emph{subtracted} to give the upper-half answer.
This add/subtract pair is the famous ``butterfly,'' and it is the reason we get $n$
output values from only $n/2$ subproblem values plus $n/2$ multiplications.
\end{intuitionbox}

\begin{center}
\begin{tikzpicture}[>=Stealth,scale=1.1,
  dot/.style={circle,fill=structurepurple,inner sep=1.5pt}]
\node[dot,label=left:{$y^{[0]}_k$}] (a) at (0,1) {};
\node[dot,label=left:{$y^{[1]}_k$}] (b) at (0,0) {};
\node[dot,label=right:{$y_k = y^{[0]}_k + \omega_n^k y^{[1]}_k$}] (c) at (3,1) {};
\node[dot,label=right:{$y_{k+n/2} = y^{[0]}_k - \omega_n^k y^{[1]}_k$}] (d) at (3,0) {};
\draw[->,sourceblue,thick] (a)--(c);
\draw[->,sourceblue,thick] (a)--(d);
\draw[->,quotientorange,thick] (b)--node[pos=0.4,above,font=\scriptsize,sloped]{$\times\,\omega_n^k$, $+$}(c);
\draw[->,quotientorange,thick] (b)--node[pos=0.4,below,font=\scriptsize,sloped]{$\times\,\omega_n^k$, $-$}(d);
\end{tikzpicture}
\end{center}

\subsection{The algorithm}

\begin{algorithm}[H]
\caption{\textsc{Recursive-FFT}$(\mathbf{a})$}
\KwIn{A coefficient vector $\mathbf{a}=(a_0,\dots,a_{n-1})$ with $n$ a power of $2$.}
\KwOut{$\mathbf{y} = \mathrm{DFT}_n(\mathbf{a})$, i.e.\ $y_k = A(\omega_n^k)$.}
\If{$n = 1$}{\Return $\mathbf{a}$\tcp*{a constant: $A(\omega_1^0)=a_0$}}
$\omega_n \gets e^{2\pi i / n}$;\quad $\omega \gets 1$\;
$\mathbf{a}^{[0]} \gets (a_0, a_2, \dots, a_{n-2})$\;
$\mathbf{a}^{[1]} \gets (a_1, a_3, \dots, a_{n-1})$\;
$\mathbf{y}^{[0]} \gets \textsc{Recursive-FFT}(\mathbf{a}^{[0]})$\;
$\mathbf{y}^{[1]} \gets \textsc{Recursive-FFT}(\mathbf{a}^{[1]})$\;
\For{$k \gets 0$ \KwTo $n/2 - 1$}{
  $t \gets \omega \cdot y^{[1]}_k$\tcp*{twiddle $\times$ odd value}
  $y_k \gets y^{[0]}_k + t$\;
  $y_{k+n/2} \gets y^{[0]}_k - t$\;
  $\omega \gets \omega \cdot \omega_n$\tcp*{advance to next root: $\omega = \omega_n^{k}$}
}
\Return $\mathbf{y}$\;
\end{algorithm}

Notice the loop maintains $\omega = \omega_n^k$ as an invariant, updating it by a
single multiplication each pass rather than recomputing a power from scratch. The
base case $n=1$ returns the single coefficient unchanged, because a degree-$0$
polynomial $A(x)=a_0$ evaluates to $a_0$ everywhere, in particular at $\omega_1^0=1$.

\subsection{Correctness, gently}

\begin{formalbox}
\textbf{Claim.} \textsc{Recursive-FFT}$(\mathbf{a})$ returns $y_k = A(\omega_n^k)$
for every $k$.

\textbf{Induction on $n$ (a power of $2$).}
\emph{Base case} $n=1$: the only root is $\omega_1^0=1$, and $A(1)=a_0$, which is
what we return. \emph{Inductive step:} assume the two recursive calls are correct
on inputs of size $n/2$, so $y^{[0]}_k = A^{[0]}(\omega_{n/2}^k)$ and
$y^{[1]}_k = A^{[1]}(\omega_{n/2}^k)$. The combine step computes, for
$0\le k<n/2$,
$$
y_k = y^{[0]}_k + \omega_n^k y^{[1]}_k,\qquad
y_{k+n/2} = y^{[0]}_k - \omega_n^k y^{[1]}_k.
$$
We showed above, using the identity $A(x)=A^{[0]}(x^2)+x A^{[1]}(x^2)$ and the
halving/cancellation lemmas, that these equal $A(\omega_n^k)$ and
$A(\omega_n^{k+n/2})$ respectively. As $k$ ranges over $0,\dots,n/2-1$, the indices
$k$ and $k+n/2$ cover all of $0,\dots,n-1$. So every output is correct. By
induction the algorithm is correct for all powers of two. $\qed$
\end{formalbox}

\subsection{Running time}

Each call does the two recursive calls on size $n/2$, plus a loop of $n/2$
iterations doing constant work each, plus building the even/odd lists in $\Th(n)$.
So
$$
T(n) = 2\,T(n/2) + \Th(n).
$$
You have met this recurrence before --- it is merge sort's recurrence. By the master
theorem (case 2, since $n^{\log_2 2} = n$ matches the $\Th(n)$ work), or by drawing
the recursion tree:

\begin{center}
\begin{tikzpicture}[level distance=9mm, >=Stealth,
  every node/.style={font=\scriptsize}]
\node {$cn$}
  child {node {$c\,\tfrac n2$}
    child {node {$c\,\tfrac n4$} child{node{$\cdots$}} child{node{$\cdots$}}}
    child {node {$c\,\tfrac n4$} child{node{$\cdots$}} child{node{$\cdots$}}}}
  child {node {$c\,\tfrac n2$}
    child {node {$c\,\tfrac n4$} child{node{$\cdots$}} child{node{$\cdots$}}}
    child {node {$c\,\tfrac n4$} child{node{$\cdots$}} child{node{$\cdots$}}}};
\node[align=left,font=\scriptsize] at (4.2,0)    {row sum $cn$};
\node[align=left,font=\scriptsize] at (4.2,-0.9) {row sum $cn$};
\node[align=left,font=\scriptsize] at (4.2,-1.8) {row sum $cn$};
\node[align=left,font=\scriptsize] at (4.2,-2.7) {$\vdots\ \ (\log_2 n + 1$ rows$)$};
\end{tikzpicture}
\end{center}

Every level of the tree does a total of $cn$ work --- the subproblems shrink but
there are proportionally more of them --- and there are $\log_2 n + 1$ levels, so
$$
T(n) = \Th(n\log n).
$$
We have evaluated a polynomial at $n$ points in $n\log n$ time. That is the
breakthrough.

\begin{complexitybox}
\begin{tabular}{lll}
\toprule
Task & Coefficient form & with the FFT \\
\midrule
Evaluate at $n$ roots of unity (DFT) & $\Th(n^2)$ & $\Th(n\log n)$ \\
Interpolate from $n$ roots (inverse DFT) & $\Th(n^2)$ & $\Th(n\log n)$ \\
Multiply two degree-$n$ polynomials & $\Th(n^2)$ & $\Th(n\log n)$ \\
Pointwise multiply $n$ values & $\Th(n)$ & $\Th(n)$ \\
\bottomrule
\end{tabular}
\end{complexitybox}

\section{Coming back: the inverse FFT}

We can travel coefficient $\to$ point-value fast. To finish the multiplication
pipeline we need the return trip, point-value $\to$ coefficient, also fast. The
miracle is that the inverse DFT is \emph{almost the same algorithm}.

Write the DFT as a matrix-vector product $\mathbf{y} = V_n\,\mathbf{a}$, where
$V_n$ is the Vandermonde matrix at the roots of unity: its $(k,j)$ entry is
$\omega_n^{kj}$. Inverting the DFT means solving $\mathbf{a} = V_n^{-1}\mathbf{y}$.
And here the summation lemma hands us the inverse on a silver platter.

\begin{theorem}[Inverse DFT]\label{ch29:idft-thm}
The inverse of $V_n$ has $(k,j)$ entry $\tfrac1n\,\omega_n^{-kj}$. Equivalently,
$$
a_j = \frac{1}{n}\sum_{k=0}^{n-1} y_k\,\omega_n^{-kj},\qquad j=0,\dots,n-1.
$$
\end{theorem}

\begin{formalbox}
\textbf{Why.} Multiply the candidate inverse $W$ (with $W_{jk}=\tfrac1n\omega_n^{-kj}$)
by $V_n$ and look at the $(j,j')$ entry of the product:
$$
[WV_n]_{jj'} = \sum_{k=0}^{n-1} \frac1n\,\omega_n^{-kj}\,\omega_n^{kj'}
            = \frac1n\sum_{k=0}^{n-1} \omega_n^{k(j'-j)}.
$$
By the summation lemma (Lemma~\ref{ch29:summation}) the inner sum is $n$ when
$j=j'$ and $0$ otherwise. So $[WV_n]_{jj'}$ is $1$ on the diagonal and $0$ off it:
$WV_n = I$. Hence $W = V_n^{-1}$. $\qed$
\end{formalbox}

Stare at the inverse formula next to the forward one:
$$
y_k = \sum_{j} a_j\,\omega_n^{kj}
\qquad\text{vs.}\qquad
a_j = \frac1n\sum_{k} y_k\,\omega_n^{-kj}.
$$
They differ in exactly two cosmetic ways: replace $\omega_n$ by $\omega_n^{-1}$, and
divide the final answer by $n$. So we can compute the inverse DFT with the
\emph{same} \textsc{Recursive-FFT} code: feed it $\mathbf{y}$, use $\omega_n^{-1}$ in
place of $\omega_n$, and divide every output by $n$. The cost is therefore the same
$\Th(n\log n)$.

\begin{walkbox}
Take a walk and convince yourself that $\omega_n^{-1}$ is itself a primitive $n$th
root of unity --- it is just $\omega_n$ traversed clockwise instead of
counterclockwise. Then ask: does the halving lemma still hold for $\omega_n^{-1}$?
(It must, or the same code could not work.) Why? Hint: the roots of unity as a
\emph{set} do not care which direction you number them.
\end{walkbox}

\section{The whole pipeline: fast polynomial multiplication}

We can finally assemble the algorithm we promised at the very start.

\begin{algorithm}[H]
\caption{\textsc{Poly-Multiply}$(\mathbf{a},\mathbf{b})$}
\KwIn{Coefficient vectors of $A,B$, each of degree bound $n$.}
\KwOut{Coefficient vector of $C = A\cdot B$, degree bound $2n-1$.}
Let $m$ be the smallest power of $2$ with $m \ge 2n$\;
Pad $\mathbf{a}$ and $\mathbf{b}$ with zeros up to length $m$\;
$\mathbf{y}_A \gets \textsc{Recursive-FFT}(\mathbf{a})$\tcp*{evaluate $A$ at $m$ roots}
$\mathbf{y}_B \gets \textsc{Recursive-FFT}(\mathbf{b})$\tcp*{evaluate $B$ at $m$ roots}
\For{$k \gets 0$ \KwTo $m-1$}{
  $\mathbf{y}_C[k] \gets \mathbf{y}_A[k]\cdot \mathbf{y}_B[k]$\tcp*{pointwise multiply}
}
$\mathbf{c} \gets \textsc{Inverse-FFT}(\mathbf{y}_C)$\tcp*{interpolate back}
\Return $\mathbf{c}$ (real parts, rounded if inputs were integers)\;
\end{algorithm}

\begin{focusbox}
Three FFTs ($\Th(m\log m)$), one pointwise loop ($\Th(m)$), and $m = \Th(n)$.
Total: $\Th(n\log n)$. We multiply two polynomials --- equivalently, convolve two
sequences --- in $n\log n$ time. The schoolbook $n^2$ is broken.
\end{focusbox}

\subsection{A complete numerical trace}\label{ch29:lab}

Let us multiply $A(x) = 1 + 2x$ and $B(x) = 3 + 4x$. By hand,
$C(x) = 3 + 10x + 8x^2$, so we are hoping to recover coefficients $(3,10,8)$. The
product has degree bound $3$, so we pad to $m=4$ and use the $4$th roots of unity
$$
\omega_4^0 = 1,\quad \omega_4^1 = i,\quad \omega_4^2 = -1,\quad \omega_4^3 = -i.
$$
Padded coefficients: $\mathbf{a}=(1,2,0,0)$, $\mathbf{b}=(3,4,0,0)$.

\textbf{Step 1: evaluate $A$.} $A(x) = 1 + 2x$ (the padding zeros do nothing), so
$$
y_A = \big(A(1),A(i),A(-1),A(-i)\big) = (3,\; 1+2i,\; -1,\; 1-2i).
$$
\textbf{Step 1b: evaluate $B$.} $B(x)=3+4x$, so
$$
y_B = (7,\; 3+4i,\; -1,\; 3-4i).
$$
\textbf{Step 2: pointwise multiply.}
$$
y_C = \big(3\cdot 7,\ (1+2i)(3+4i),\ (-1)(-1),\ (1-2i)(3-4i)\big).
$$
Work out the middle two: $(1+2i)(3+4i) = 3 + 4i + 6i + 8i^2 = 3 + 10i - 8 = -5 + 10i$,
and its conjugate $(1-2i)(3-4i) = -5 - 10i$. So
$$
y_C = (21,\ -5+10i,\ 1,\ -5-10i).
$$
\textbf{Step 3: inverse DFT.} Use $a_j = \tfrac14\sum_k y_{C,k}\,\omega_4^{-kj}$ with
$\omega_4^{-1} = -i$. For example, $c_0 = \tfrac14(21 + (-5+10i) + 1 + (-5-10i))
= \tfrac14(12) = 3.$ And
$$
c_1 = \tfrac14\big(21\cdot 1 + (-5+10i)(-i) + 1\cdot(-1) + (-5-10i)(i)\big).
$$
Compute the pieces: $(-5+10i)(-i) = 5i + 10 = 10 + 5i$; \ $(-5-10i)(i) = -5i + 10 =
10 - 5i$. Sum: $21 + (10+5i) - 1 + (10-5i) = 40$, so $c_1 = 40/4 = 10$. A similar
computation gives $c_2 = 8$ and $c_3 = 0$.

\begin{labbox}
\textbf{Trace it yourself.} You just watched $(1+2x)(3+4x)$ come out as
$3 + 10x + 8x^2$ through the FFT machinery. Now do
$A(x)=1+x+x^2$ times $B(x)=1-x$ by the same recipe.
\begin{enumerate}
\item The product has degree $3$, so pad to $m=4$. Write the padded vectors.
\item Evaluate both at $1,i,-1,-i$. (You already know $A(i)=1+i+i^2 = i$.)
\item Multiply the four value-pairs.
\item Run the inverse formula $c_j=\tfrac14\sum_k y_{C,k}(-i)^{kj}$.
\item Check against the schoolbook answer
$(1+x+x^2)(1-x) = 1 - x^3$, i.e.\ coefficients $(1,0,0,-1)$. Did the
imaginary parts cancel to (almost) zero? They must, since real inputs give a real
product.
\end{enumerate}
\end{labbox}

\begin{trapbox}
\textbf{Five classic mistakes.}
\begin{itemize}
\item \textbf{Forgetting to pad.} If you sample $A$ and $B$ at only $n$ points, you
get $n$ values of $C$, but $C$ has degree bound $2n-1$ and needs $2n$ values. Your
interpolation will silently produce the wrong (``wrapped-around'') answer. Always
pad to length $\ge 2n$, rounded up to a power of two.
\item \textbf{Using $+\omega_n$ instead of $\omega_n^{-1}$ on the way back}, or
\textbf{forgetting the $1/n$}. Either one corrupts the inverse transform.
\item \textbf{Assuming the imaginary parts are exactly zero.} With floating-point
arithmetic the recovered coefficients carry tiny imaginary residue like
$10^{-15}i$; take real parts and (for integer problems) round.
\item \textbf{Thinking the FFT needs $n$ prime, or $n$ anything other than a power of
two.} The radix-$2$ FFT in this chapter needs $n$ to be a power of two; pad up to
one. (More general FFTs exist, but start here.)
\item \textbf{Confusing degree with degree bound.} A polynomial of \emph{degree}
$d$ has \emph{degree bound} $d+1$, i.e.\ $d+1$ coefficients. Off-by-one errors here
cause the worst padding bugs.
\end{itemize}
\end{trapbox}

\section{Why the same code runs forward and backward}

It is worth one more pass on the elegant symmetry, because it is easy to use the
inverse FFT mechanically without seeing \emph{why} the universe is so kind.

\begin{intuitionbox}
The forward DFT mixes the coefficients into values by ``winding'' them around the
circle counterclockwise (multiplying by $\omega_n = e^{+2\pi i/n}$). The inverse
unwinds them by going clockwise ($\omega_n^{-1} = e^{-2\pi i/n}$) and then dividing
out the factor $n$ that the winding introduced. Because both directions use roots of
unity --- and the roots of unity are the \emph{same set} whether you number them
clockwise or counter --- the halving and butterfly structure is identical. The FFT
is a machine that runs equally well in reverse; you only change the gear ratio
($\omega \to \omega^{-1}$) and a final scale ($\div n$).
\end{intuitionbox}

\begin{aestheticbox}
This forward/backward symmetry is exactly what makes the FFT the workhorse of signal
processing. To filter a sound (remove hiss, boost bass, cancel an echo) you FFT into
the frequency domain, tweak the frequency values, and inverse-FFT back to a waveform
--- and both legs cost $n\log n$. Convolution in the time domain becomes
\emph{pointwise multiplication} in the frequency domain, which is the same fact we
used to multiply polynomials. Audio plug-ins, noise-cancelling headphones, MRI
reconstruction, and the modulation scheme in every Wi-Fi router are all standing on
this one $n\log n$ algorithm.
\end{aestheticbox}

\begin{historybox}
The divide-and-conquer FFT is usually credited to James Cooley and John Tukey
(1965), who were partly motivated by detecting Soviet nuclear tests from
seismic data. But the core idea is far older: Carl Friedrich Gauss used essentially
the same trick around 1805 to interpolate the orbits of asteroids, before Fourier's
work was even published --- his notes lay buried in his collected works in
non-standard notation for over a century. This material is CLRS Chapter~30
(``Polynomials and the FFT''): Section~30.1 covers the two representations and the
multiplication strategy, Section~30.2 develops the DFT and \textsc{Recursive-FFT},
and Section~30.3 treats efficient implementations (the iterative, in-place FFT with
bit-reversal). For the number theory behind doing FFTs in modular arithmetic --- the
``Number-Theoretic Transform,'' which avoids floating-point entirely --- see the
next chapter on number-theoretic algorithms.
\end{historybox}

\section{A note on making it iterative and fast in practice}

The recursive code is the clearest way to \emph{understand} the FFT, but production
implementations are usually \emph{iterative} and operate in place. Two ideas drive
this.

First, \emph{bit-reversal}. If you trace which input coefficient ends up feeding each
leaf of the recursion, you find that coefficient $a_j$ lands in the slot whose index
is the binary digits of $j$ \emph{reversed}. For $n=8$, $a_1$ ($001_2$) lands at
position $100_2 = 4$, $a_3$ ($011_2$) lands at $110_2 = 6$, and so on. Permute the
input by bit-reversal once at the start, and the recursion's even/odd splitting is
already laid out for you.

Second, \emph{building up} instead of down. After bit-reversal, combine adjacent
pairs with butterflies into transforms of size $2$, then size $4$, then size $8$,
$\dots$ up to size $n$ --- $\log_2 n$ passes, each touching all $n$ elements with
$n/2$ butterflies. No recursion stack, no array allocation per call, excellent
cache behavior. Same $\Th(n\log n)$, much smaller constant.

The picture below shows the shape of the iterative FFT for $n=8$. The inputs enter on
the left in bit-reversed order; three columns of butterflies (combining blocks of
size $2$, then $4$, then $8$) carry them rightward into the finished transform
$y_0,\dots,y_7$. Each ``$\bowtie$'' is one butterfly: an add and a subtract sharing a
twiddle multiply, exactly as in the recursive combine step.

\begin{center}
\begin{tikzpicture}[>=Stealth,xscale=1.55,yscale=0.62,
  lbl/.style={font=\scriptsize,inner sep=1pt},
  bf/.style={structurepurple,thick}]
% rows 0..7 from top(0) to bottom(7); inputs in bit-reversed order
\foreach \r/\j in {0/0,1/4,2/2,3/6,4/1,5/5,6/3,7/7}{
  \node[lbl,left] at (-0.15,-\r) {$a_{\j}$};
  \node[lbl,right] at (3.15,-\r) {$y_{\r}$};
  \draw[gray!55] (0,-\r)--(3,-\r);
}
% stage column x positions: 1,2,3 are right edges; draw butterfly cross links
% Stage 1: pairs (0,1)(2,3)(4,5)(6,7) at x in [0,1]
\foreach \p in {0,2,4,6}{
  \draw[bf] (0.0,-\p)--(1.0,{-\p-1});
  \draw[bf] (0.0,{-\p-1})--(1.0,-\p);
}
% Stage 2: blocks of 4: pairs distance 2: (0,2)(1,3)(4,6)(5,7) at x in [1,2]
\foreach \b in {0,4}{
  \foreach \o in {0,1}{
    \draw[bf] (1.0,{-\b-\o})--(2.0,{-\b-\o-2});
    \draw[bf] (1.0,{-\b-\o-2})--(2.0,{-\b-\o});
  }
}
% Stage 3: block of 8: pairs distance 4: (0,4)(1,5)(2,6)(3,7) at x in [2,3]
\foreach \o in {0,1,2,3}{
  \draw[bf] (2.0,{-\o})--(3.0,{-\o-4});
  \draw[bf] (2.0,{-\o-4})--(3.0,{-\o});
}
\node[font=\tiny,gray] at (0.5,0.7) {size-2};
\node[font=\tiny,gray] at (1.5,0.7) {size-4};
\node[font=\tiny,gray] at (2.5,0.7) {size-8};
\end{tikzpicture}
\end{center}

\begin{complexitybox}
\textbf{Recursive vs.\ iterative FFT.} Both are $\Th(n\log n)$. The iterative,
in-place version (CLRS 30.3) wins by a constant factor: it avoids per-call
allocation, computes twiddle factors with one multiply per use, and streams through
memory in a cache-friendly pattern. For the asymptotics in this book either is fine;
for real software, write the iterative one.
\end{complexitybox}

\begin{recoverybox}
If this chapter went by fast, here is the one paragraph to keep. A polynomial can be
written by its coefficients \emph{or} by its values at chosen points; multiplication
is slow ($n^2$) in coefficients but fast ($n$) in values. Converting between the two
forms is normally $n^2$, but if you choose the points to be the $n$th roots of unity
--- evenly spaced on the unit circle --- then a divide-and-conquer split on even/odd
coefficients turns the conversion into $n\log n$. That conversion is the FFT; running
it forward, multiplying pointwise, and running it backward multiplies polynomials in
$\Th(n\log n)$.
\end{recoverybox}

\section{Exercises}

\begin{exercise}\label{ch29:ex-pv}
Convert $A(x) = 2 - x + x^2$ from coefficient form to a point-value representation
using the points $x=-1,0,1$. Then, pretending you only had those three pairs,
reconstruct the polynomial with Lagrange interpolation and confirm you get the
original coefficients back.
\end{exercise}

\begin{exercise}\label{ch29:ex-conv}
Write out the convolution sum for $c_k$ when $A=(1,2,3)$ and $B=(4,5,6)$, computing
all five coefficients $c_0,\dots,c_4$ of the product by hand. How many scalar
multiplications did the schoolbook method use? Now state the count for general degree
bound $n$.
\end{exercise}

\begin{exercise}\label{ch29:ex-roots}
List the six $6$th roots of unity in the form $a+bi$ (use exact values where you can,
decimals otherwise). Verify directly that their sum is $0$, illustrating the
summation lemma.
\end{exercise}

\begin{exercise}\label{ch29:ex-halving}
Compute the squares of all eight $8$th roots of unity and confirm by hand that the
results are exactly the four $4$th roots of unity, each appearing twice (the halving
lemma).
\end{exercise}

\begin{exercise}\label{ch29:ex-trace}
Run \textsc{Recursive-FFT} on the input $\mathbf{a} = (1,2,3,4)$ ($n=4$). Show the
even/odd split, the two size-$2$ recursive results, and the butterfly combine. Your
answer should be the four values $A(1),A(i),A(-1),A(-i)$. Check by evaluating
$A(x)=1+2x+3x^2+4x^3$ directly.
\end{exercise}

\begin{exercise}\label{ch29:ex-inverse}
Take the output vector you computed in Exercise~\ref{ch29:ex-trace} and run the
inverse DFT on it (use $\omega_4^{-1}=-i$ and divide by $4$). You should recover
$(1,2,3,4)$. Show at least the computations for $a_0$ and $a_2$.
\end{exercise}

\begin{exercise}\label{ch29:ex-multiply}
Use the full \textsc{Poly-Multiply} pipeline to compute $(1 + x + x^2)(2 + 3x)$.
Pad to the right power of two, do the two forward FFTs, multiply pointwise, and
inverse-transform. Confirm the coefficients match the schoolbook product.
\end{exercise}

\begin{exercise}\label{ch29:ex-bitrev}
For $n=8$, write the bit-reversal permutation: for each index $0,\dots,7$ give its
$3$-bit binary, reverse the bits, and report the destination index. Which indices are
fixed points (map to themselves)?
\end{exercise}

\begin{exercise}\label{ch29:ex-recurrence}
The recurrence for \textsc{Recursive-FFT} is $T(n) = 2T(n/2) + cn$. Solve it two
ways: (a) by drawing the recursion tree and summing the per-level work; (b) by the
master theorem, identifying which case applies and why. Confirm both give
$\Th(n\log n)$.
\end{exercise}

\begin{exercise}\label{ch29:ex-bignum}
\emph{(Challenge.)} Explain how to multiply two $n$-digit integers in $\Th(n\log n)$
time using the FFT. (Hint: an integer $\overline{d_{n-1}\cdots d_1 d_0}$ in base $10$
is the polynomial $\sum d_j\,10^j$ evaluated at $x=10$. Multiply the polynomials,
then handle carries.) What goes wrong if you forget to pad, and what is the role of
carrying after the convolution?
\end{exercise}

\begin{exercise}\label{ch29:ex-strang}
\emph{(Reflection.)} The summation lemma says the roots of unity sum to zero for any
$k$ not a multiple of $n$. Explain geometrically why this must be true by thinking of
the roots as vectors pointing from the origin, equally spaced around the circle.
(What is the ``center of mass'' of evenly spaced points on a circle?)
\end{exercise}
